\section*{Chapitre \ref{ch:m3:crude-oil} --- Pétrole brut}

\begin{solution}{exo:m3:crude-oil:1}
$141.5/0.87 - 131.5 = 31.1$° API~: plus lourd (plus dense) qu'un brut à 39.6°.
\end{solution}

\begin{solution}{exo:m3:crude-oil:2}
Le 25 novembre 2026 est un jour ouvré, donc la négociation prend fin trois jours ouvrés
avant cette date~: vendredi 20 novembre 2026.
\end{solution}

\begin{solution}{exo:m3:crude-oil:3}
La production d'un seul champ baisse avec le temps, et une petite offre livrable peut
faire l'objet d'un squeeze~: celui qui contrôle la plupart des cargaisons d'un mois peut faire bouger le prix. Plusieurs qualités,
avec les ajustements de qualité et la règle de la qualité la plus compétitive, rendent l'offre
livrable plus large et l'indice de référence plus difficile à manipuler.
\end{solution}

\begin{solution}{exo:m3:crude-oil:4}
$80 + 0.04 \times (36 - 38) - 0.12 \times (0.60 - 0.40)/0.1 = 80 - 0.08 - 0.24 =
\qty{79.68}{\$/\barrel}$.
\end{solution}

\begin{solution}{exo:m3:crude-oil:5}
$25 - 0.40 = \qty{24.60}{\$/\barrel}$. Un prix de l'échéance la plus proche de 18~\$ implique un coût de stockage de
7~\$ par baril-mois~: le stockage à 0.40~\$ n'est pas disponible pour ceux qui pourraient arbitrer,
ou bien le contrat à l'échéance la plus proche est liquidé plus vite qu'on ne trouve de capacité de stockage.
\end{solution}

\begin{solution}{exo:m3:crude-oil:6}
$(71.35 + 2.20) \times 2\,000\,000 = $ 147.1~millions de dollars USD.
\end{solution}

\begin{solution}{exo:m3:crude-oil:7}
25 jours au-dessus de 5~\$~; la fourchette la plus large est de 58.06~\$ (20 avril 2020), la
deuxième est de 8.49~\$.
\end{solution}

\begin{solution}{exo:m3:crude-oil:8}
Un acheteur d'un contrat à livraison physique reçoit du pétrole en un lieu et à une date précis.
Si l'on n'a nulle part où le stocker et que le déplacer coûte de l'argent, le pétrole à ce lieu et
à cette date a une valeur négative pour quiconque n'a pas de stockage~: c'est un passif dont il faut
se débarrasser. «~Le pétrole gratuit~» vaut sa valeur ailleurs et plus tard, moins le coût du
stockage et du transport, coût qui peut dépasser cette valeur.
\end{solution}

\begin{solution}{pb:m3:crude-oil:1}
\textbf{1.} 1\,000\,000 barils. \textbf{2.} 21 avril 2020~: le 25 avril était un
samedi, le dernier jour ouvré avant celui-ci était le vendredi 24 avril, et trois jours ouvrés
avant celui-ci, c'est le mardi 21 avril~; le 20 avril est la veille. \textbf{3.} 18.27
millions de dollars USD. \textbf{4.} 6.76~\$, un coût de stockage de 6.76~\$ par baril-mois. \textbf{5.}
Non~: environ 17 fois le loyer normal~; le stockage était déjà rare.
\textbf{6.} $(-37.63 - 18.27) \times 1\,000\,000 = -$ 55.90~millions de dollars USD.
\textbf{7.} 58.06~\$ par baril-mois. \textbf{8.} Environ 145 fois. \textbf{9.}
Personne parmi ceux qui pouvaient encore prendre livraison et stocker n'avait plus de place, si bien que l'arbitrage
de la \cref{prop:m3:crude-oil:floor} ne pouvait pas être mené à grande échelle. \textbf{10.}
$(58.06 - 0.40) \times 1\,000\,000 = $ 57.66~millions de dollars USD.
\textbf{11.} 6.76~\$ le baril, 6.76~millions de dollars USD. \textbf{12.} 6.76~\$ le baril, 27~\%
du prix de juin, en un mois. \textbf{13.} Les échéances ultérieures étaient plus éloignées
de la livraison et valorisées avec la reprise attendue~; un fonds y détenant des positions n'était pas
contraint de vendre pendant la fenêtre de règlement-livraison. \textbf{14.} Le prix a été fixé par la moyenne des transactions sur spreads
entre 14:28 et 14:30, moment où les détenteurs sans stockage étaient forcés de vendre~: une petite
fenêtre au pire moment. \textbf{15.} Le modèle de Bachelier, qui admet un sous-jacent
négatif~; les modèles lognormaux ne peuvent pas valoriser des options sur un contrat à terme négatif.
\textbf{16.} Les détenteurs de capacités de stockage ou de pipeline, les raffineurs et les traders en position vendeuse sur le
contrat de mai disposant d'une capacité de livraison~: pour eux, le pétrole à $-\$37.63$ revenait à être payé pour le stocker.
\textbf{17.} Un indice de référence à livraison physique valorise la marchandise en un lieu et
à une date donnée, y compris le coût de sa prise en charge~; un contrat réglé en espèces valorise un indice et ne peut
pas être contraint à la livraison. \textbf{18.} Non~: l'ICE Brent est réglé en espèces sur un indice d'un
marché maritime où les navires servent de stockage~; le prix spot du Brent de l'EIA n'est jamais descendu sous
9.12~\$ (21 avril 2020). \textbf{19.} \emph{Résultat nommé~:} 58.06~\$ par baril pour un
mois, environ 145 fois le loyer illustratif. \textbf{20.} De la marchandise au
lieu et à la date de livraison, net du coût de prise en charge et de conservation sur place.
\end{solution}

\begin{solution}{iq:m3:crude-oil:1}
Le \omterm{def:m3:crude-oil:brent}{Dated Brent} est l'évaluation par une agence des cargaisons physiques de brut de la mer du Nord en chargement dans les
semaines à venir~; les contrats à terme sur l'ICE Brent sont négociés en bourse, à plusieurs mois d'échéance, et
réglés en espèces sur un indice du marché à terme. Les CFD et le marché
à terme les relient.
\iqlookfor{physique contre financier, la dimension temporelle, les instruments de liaison.}
\end{solution}

\begin{solution}{iq:m3:crude-oil:2}
Lorsque le prix porte sur une livraison en un lieu et à un moment où prendre possession de la marchandise coûte
plus que sa valeur~: stockage saturé, absence de transport, coûts d'élimination (électricité lors d'une
heure de congestion, gaz dans un réseau plein). Le financier suit le physique à la livraison.
\iqlookfor{le prix porte sur un bien livré~; le stockage et la logistique ne fixent un plancher que s'ils sont disponibles.}
\end{solution}

\begin{solution}{iq:m3:crude-oil:3}
Savoir si la série continue a été rollée avant l'avant-dernier jour, comme le ferait tout fonds réel,
si la stratégie pouvait même détenir le contrat de mai (livraison,
seuils de déclaration), si l'exécution au règlement-livraison était réalisable, et
si le P\&L provient d'une seule journée~; une série de rendements logarithmiques casse sur un prix négatif.
\iqlookfor{conventions de roll, faisabilité, le prix négatif comme piège d'erreur de données.}
\end{solution}

\begin{solution}{iq:m3:crude-oil:4}
Un spread de 1.50~\$ face à 0.40~\$ de stockage rapporte 1.10~\$ en cash-and-carry~: acheter le contrat
à l'échéance la plus proche, prendre livraison, stocker, vendre le second mois. Risques~: stockage et pipeline
planification, financement, qualité et pertes, et la possibilité que le spread s'élargisse
encore avant la livraison, ce qui coûte en appels de marge.
\iqlookfor{le carry trade, et ses risques opérationnels et de liquidité.}
\end{solution}

\begin{solution}{iq:m3:crude-oil:5}
$\ln F$ n'est pas défini pour $F \le 0$~: les prix, les volatilités implicites et les grecques échouent.
Passer à un modèle normal (de Bachelier) ou à une lognormale décalée, revaloriser la surface dans
ce modèle, et faire en sorte que le code refuse un logarithme d'un prix non positif plutôt que de renvoyer
NaN.
\iqlookfor{Bachelier, modèles décalés, code défensif.}
\end{solution}

\begin{solution}{iq:m3:crude-oil:6}
Plusieurs qualités signifient davantage d'offre livrable et davantage de vendeurs, si bien qu'accaparer
la qualité la moins chère devient coûteux~; le prix est celui de la qualité la plus compétitive. Le coût tient
à la complexité~: les primes de qualité doivent être fixées et révisées, et la qualité de référence peut
changer d'un jour à l'autre, ce qui ajoute de la base pour les hedgers.
\iqlookfor{offre livrable face aux squeezes~; coûts d'ajustement de qualité et de base.}
\end{solution}
